\section{Coverage rechecked, and a stronger use of multiplicities}
\label{sec:kt-coverage}

\subsection{The source classes and normalization}
The source commit is
\texttt{7cbf7a0c32fca1badf05931879a37c3f8a2b03b4}.
Paths in this paragraph are relative to
\texttt{src/Development/AutomorphicArithmetic/AutomorphicVarieties/}.
The exact predicate is
\texttt{ModularityArgument/ThresholdArguments/DenseThresholdCases/NumberU3US.lean},
lines 27--29: \texttt{Adm3U\_s} extends \texttt{Adm3Alt} and adds
\texttt{Th.Unl}. The latter base predicate, in \texttt{AfStmt.lean},
lines 29--32, supplies irreducibility, the automorphic condition for
$r=\rho\chi_\ell^{-1}$, and
$\operatorname{Level}(\ell,2,(e,0,-e),\rho)$.
\texttt{ModularityArgument/AutomorphyLifting/IsAutCsdA.lean},
lines 11--12, retains both \texttt{AutCsd0} and the implication
from continuity to \texttt{WitA}.
In the parallel \texttt{Foundations/GaloisRepresentations/AdmissibleResidues/}
directory, \texttt{IsAutomorphicOverImagQuad.lean}, lines 35--36,
defines \texttt{AutCsd0} by a regular algebraic conjugate-self-dual
\emph{cuspidal} $\mathrm{GL}_3$ representation with attachment.

\texttt{AutomorphicForms/BlochKatoSelmerVanishes.lean},
lines 66--67 and 89--114, defines the root-of-unity and local
level conditions. Because $\mu_2\subset\mathbf Q$, every
\texttt{KillsMu 2} character is trivial; the same level witnesses
therefore work at modulus one.
\texttt{BuildingsResidueChars/LevelTwistCyclotomicPow.lean},
lines 19--61, changes the exponent vector to
$(e-1,-1,-e-1)$ after the common twist.
Continuity is part of \texttt{Level}, so \texttt{WitA} is available.
\texttt{ModularityArgument/AutomorphyLifting/AdmissibleLevelRestriction/UnlAllOf.lean},
lines 36--38, gives the required above-$\ell$ nonadjacent flag;
the common twist preserves its inequalities.
\texttt{KeepLevel.lean}, lines 30--82 and 95--101, then supplies
an everywhere-finite-unramified cusp representation, its infinitesimal
character $(e,0,-e)$, and Frobenius--Satake compatibility at every
place away from $\ell$.
At modulus one the group \texttt{KLev} is the full integral adelic
maximal compact: this follows directly from
\texttt{AutomorphicForms/KLevTop.lean}, \texttt{mem\_KLev},
lines 17--48. All these relevant proofs and definitions were read
directly in the pinned files; their imported foundations were not
independently rebuilt in Lean.

The rank-two predicate is
\texttt{ModularityArgument/AutomorphyLifting/AutomorphyDescent/AutNewformW.lean},
\texttt{Adm2A}, lines 14--20.
It supplies an irreducible continuous $W$, a weight-$2e+1$
newform attached to $W\chi_\ell^{-e}$, the actual
$\operatorname{Level}(\ell,6,(e,-e),W)$, and
$\rho\simeq1\oplus W|_{G_E}$, as well as its extra
\texttt{WitAQ} and determinant conditions.
The definition of \texttt{Newform} and its attachment in
\texttt{AutomorphicForms/Newform.lean}, lines 21--33 and 43--47,
includes normalization, a minimal-level clause, semisimplicity,
and cofinite equality of full Frobenius characteristic polynomials.

\subsection{Published input for the compact transfer}

\paragraph{Rechecked inherited theorem: K17-U3-COVERAGE and
K17-U3-GENERAL-COVERAGE.}
For every even $e\geq2$ and odd $\ell$, every source
$\operatorname{Adm3U\_s}(\iota,e,2)$ representation is covered
by the full compact space $M_{e-1}$, and its split-prime trace is
$\alpha/p^e$, where $\alpha$ is an eigenvalue of
$p^{e-1}T_p$ on that space.
The following checks explain why the inherited proof uses the
published theorems within their stated scope.

KMSW, \emph{Endoscopic classification of representations: inner forms
of unitary groups}, arXiv:1409.3731v3, Section~0.3.3, pp.~17--18,
identifies the Hermitian-space unitary groups as pure inner forms;
the standard positive form is one of them, and in odd rank its finite
local groups are quasi-split.
Propositions~1.3.1 and 1.3.3, pp.~67 and 70--71, attach and localize
the simple generic parameter of the cusp representation.
The parity criterion of Section~1.2.3 and Lemma~1.2.5, pp.~59--61,
applies because the distinct infinity characters
$(z/\bar z)^e,1,(z/\bar z)^{-e}$ all have positive conjugate parity.
The reduced global component group is trivial by the definitions
and formula (1.3.6), Section~1.3.4, pp.~69--70.
The compact member at infinity has highest weight $(e-1,0,1-e)$.
The real-packet relevance assertion is Bellaiche--Chenevier,
\emph{Families of Galois representations and Selmer groups},
Appendix~A.11, Remark~A.11.8 and its footnote~142, p.~236;
the entire remark and its explicit local-parameter argument were
read here. Only this local statement is used, not that appendix's
conditional global-parameter formalism.

The global multiplicity formula used is precisely KMSW
Theorem~5.0.5 and its preceding paragraph, p.~206.
That paragraph verifies its two hypotheses for generic parameters
and pure inner twists. Its proof, read in full here, gives the
direct sum over the packet selected by the global character.
Section~1.7, pp.~88--89, shows that the product of the local central
characters descends to the reduced component group. Since that
group is trivial here, independently chosen relevant local members
have multiplicity one.
No unrestricted version of the starred Theorem~1.7.1 is assumed.

Shin, \emph{Galois representations arising from some compact Shimura
varieties}, Annals of Mathematics 173 (2011), Theorem~1.2 and
Corollary~1.3, pp.~1647--1648, applies to regular algebraic
conjugate-self-dual cuspidal representations of odd rank.
It makes every finite local parameter tempered; its additional
slight-regularity condition applies only in even rank.
Corollary~7.9 and its printed proof, p.~1728, were also read.
At a split prime the required vector is the ordinary
$\mathrm{GL}_3(O_{E_w})$-fixed vector.
At an inert prime the parameter is unramified and the hyperspecial
member exists by KMSW Proposition~1.5.2(5), p.~79.

At the ramified prime three a full-stabilizer check is necessary.
The central character of the global $\Pi$ is unramified at every
finite place and trivial at infinity, hence trivial because $E$
has class number one. The roots of its local unramified parameter
are consequently $\{a,1,a^{-1}\}$, with $|a|=1$.
For $F=\mathbf Q_3$ and $L=E_3$, choose the unramified character
$\chi:L^\times\to\mathbf C^\times$ with Frobenius value $a$.
On the diagonal torus of $U(2,1)$ take the character
$\lambda(\operatorname{diag}(\alpha,\beta,\bar\alpha^{-1}))
=\chi(\alpha)$.
Its stable torus lift is $(\chi,1,\chi^{-1})$.
The induced $\mathrm{GL}_3(L)$ representation is irreducible
because these unramified characters are unitary.

Adler--Lansky, \emph{Depth-zero base change for ramified $U(2,1)$},
arXiv:0807.1528v2, Section~4.4, pp.~14--15, lists the reducibility
possibilities. The nonunitary case is excluded, and the case whose
restriction to $F^\times$ is the ramified quadratic character is
excluded because $\chi$ is unramified. The remaining reducible
case is the nontrivial unramified quadratic $\chi$, whose two
constituents form one packet.
Proposition~6.1(1),(4), p.~21, and Proposition~7.1(i),(iv), p.~28,
identify the stable base change of the irreducible or two-member
packet as the required $\Pi_L$.
The proof of Proposition~7.1 refers to the earlier unramified paper:
arXiv:math/0601695, Proposition~5.1, pp.~19--20 (published
numbering Proposition~4.1). Its entire proof, read here, derives
the used cases directly from Rogawski, Proposition~4.10.2 and
the paragraph before Theorem~13.2.1.
Thus neither the separate supercuspidal restriction
$q>3$, $q\ne7$, nor the depth-preservation Hypothesis~1.5 is an
input to these principal-series cases.

For completeness the standard lattice is the full vertex-$y$
stabilizer in Adler--Lansky Sections~4.1 and 4.2.1, pp.~10--12.
Choose $b\in\mathbf Z_3$ with $b^2=-2$ by Hensel's lemma.
The vector $(1,1,b)$ is primitive isotropic; pairing it with an
integral vector of pairing one and subtracting half the latter's
norm gives an integral hyperbolic plane. If $d$ is the norm of
the complementary unit line and $B$ is the basis matrix,
$-d=N(\det B)$; scaling the last vector by $(\det B)^{-1}$
makes its norm $-1$. This identifies the form with their matrix
$\left(\begin{smallmatrix}0&0&1\\0&-1&0\\1&0&0\end{smallmatrix}\right)$.
Their full reductive quotient is $O_3(k)$, including its central
component, and the vertex is special. Iwasawa decomposition gives
one full-stabilizer-fixed vector in the principal series because
both $\lambda$ and the modulus are trivial on $B(F)\cap K_3$.
Exactness of compact-open invariants in characteristic zero places
such a vector in some constituent of the required packet.
Global assembly by KMSW now proves the inherited coverage.

Finally \texttt{FrobEig} uses geometric Frobenius, whereas
$\chi_\ell(\sigma)=p$ for arithmetic Frobenius.
Therefore $r(\sigma^{-1})=p\rho(\sigma^{-1})$.
The source's \texttt{algSat},
\texttt{BlochKatoSelmerVanishes.lean}, lines 69--82, multiplies
the unitary Satake roots in rank three by $p$.
The two factors cancel, and Gross's minuscule normalization
$T_p=p\sum\beta$ gives the asserted $\alpha/p^e$.

\subsection{Published input for the rank-two conductor}

\paragraph{Rechecked inherited theorem: K17-GL2-CONDUCTOR.}
For $e\geq1$ and $\ell>3$, every
$\operatorname{Adm2A}(\iota,e,6)$ is represented by a normalized
newform $f\in S_{2e+1}(\Gamma_0(3),\chi_{-3})$, of primitive
conductor exactly three, and has split-prime trace
$1+a_p(f)/p^e$.

Here is a reconstructed proof with the exact source scopes.
The source attachment identifies $R=W\chi_\ell^{-e}$ with
the semisimple Galois representation of the primitive eigenform,
by Chebotarev and Brauer--Nesbitt.
Every \texttt{KillsMu 6} character factors through
$\operatorname{Gal}(\mathbf Q(\mu_3)/\mathbf Q)$, so it is
$1$ or $\chi_{-3}$.
At $q\ne\ell$ the actual inertia action is a direct sum of these
characters. It has finite image and zero monodromy.
Thus the local parameter is unramified at $q\ne3,\ell$, and at
three it is tame with these two characters and zero monodromy.
Carayol, \emph{Sur les representations $\ell$-adiques associees
aux formes modulaires de Hilbert}, Section~0.7, Theorem~A,
pp.~410--411, gives the full away-from-$\ell$ compatibility.
Its extra discrete-series hypothesis, Section~0.3, p.~409,
is imposed only when the totally real degree is even, so it
does not apply over $\mathbf Q$.

At $\ell$, the local flag after twisting has characters
$\delta_1$ and $\delta_2\chi_\ell^{-2e}$ with unramified
continuous $\delta_i$.
Their Frobenius values are $\ell$-adic units because the image
of a compact group under a continuous character is compact.
Thus the two crystalline characters have slopes differing by $2e$.
Saito, \emph{Galois representations in arithmetic geometry},
Section~2c, Theorem~3, p.~9, with the definitions on pp.~7--8,
gives potential semistability and full Frobenius-semisimplified
Weil--Deligne compatibility for classical modular forms over
$\mathbf Q$, including monodromy. This author-stated theorem
refers explicitly to Saito's 1997 \emph{Modular forms and $p$-adic
Hodge theory}; the original 1997 proof was not reread here.

After a finite extension over which $R$ is semistable, exactness
of $D_{\rm st}$ shows that its two slopes are those of the two
flag quotients. The relevant exactness statement is
Colmez--Fontaine, \emph{Construction des representations $p$-adiques
semi-stables}, author version, Section~4.1, Proposition~4.2,
p.~23. The rank-one period calculation is Section~4.2,
Proposition~4.4 and its proof, p.~24.
Section~3.3, p.~21, gives $N\varphi=\ell\varphi N$ and says
that $N$ lowers slopes by one. Since $2e\geq2$, the monodromy
operator in this two-dimensional module is zero.
The Weil--Deligne inertia has finite image and its two graded
characters are trivial; averaging over that finite group makes
its entire action trivial. Saito's theorem therefore makes $f$
unramified at $\ell$ as well.

The nebentypus is now unramified outside three and has conductor
dividing three. Its odd weight forces it to be odd, hence equal
to $\chi_{-3}$. Exactly one of the two inertia characters at
three is consequently nontrivial. Their tame Artin conductor
is one, and zero monodromy adds nothing, proving primitive
conductor three. Finally
$R(\sigma^{-1})=p^eW(\sigma^{-1})$ and attachment at the split
prime give the trace formula. All extra source predicates are
retained; this proof simply does not need to weaken them.

\subsection{Multiplicity is available when the whole systems are retained}

\begin{lemma}[KT-TRACE-MULTIPLICITY]
Fix an even $e\geq2$, $\ell>3$, $\ell\ne7$, and an embedding
$\iota:\overline{\mathbf Q}_\ell\simeq\mathbf C$.
Let $X$ be a finite set of pairwise nonisomorphic source
representations, each satisfying either
$\operatorname{Adm3U\_s}(\iota,e,2)$ or
$\operatorname{Adm2A}(\iota,e,6)$.
Let $P_{e,7}$ be the characteristic polynomial of
$7^{e-1}T_7$ on the full compact space $M_{e-1}$, and
$Q_{e,7}$ the characteristic polynomial of $T_7$ on the full
$S_{2e+1}(\Gamma_0(3),\chi_{-3})$.
Put
\[
 N_{e,7}=P_{e,7}(7^{2e}+7^e+1)\,Q_{e,7}(7^{2e}+1).
\]
With $s$ a geometric Frobenius at a selected prime above seven,
$t=7^{-e}+1+7^e$, and $v_\ell(\ell)=1$, one has
\[
 \sum_{\rho\in X}v_\ell(\operatorname{tr}\rho(s)-t)
                       \leq v_\ell(N_{e,7}).           \tag{KT-M.1}
\]
The polynomial evaluations are nonzero. Under the source
$\mathrm{GD}$ hypothesis and $\mathrm{OneCharU\_s}$,
$V_n\leq2v_\ell(N_{e,7})$ follows.
\end{lemma}
\begin{proof}
Partition $X$ into its irreducible rank-three branch and its
reducible $1\oplus W|_{G_E}$ branch. These branches are disjoint.
For each rank-three member choose its attached cusp representation
$\Pi$. If two members had the same $\Pi$, their twists would be
continuous semisimple representations with equal full Frobenius
characteristic polynomials outside a finite set, by
\texttt{AttachedSys},
\texttt{BlochKatoSelmerVanishes.lean}, lines 77--82.
Chebotarev and Brauer--Nesbitt would make the two members
isomorphic, contrary to the hypothesis.
Thus the chosen $\Pi$ are distinct.
The generic multiplicity formula puts their compact transfers
in distinct automorphic parameter summands. Selecting one
$K$-fixed Hecke eigenline in each summand gives linearly
independent eigenvectors in $M_{e-1}$.
Accordingly their eigenvalues at seven occur in $P_{e,7}$ with
at least the required multiplicities, even when several values
at seven coincide.

For the second branch choose the normalized primitive newform
supplied by the conductor theorem. A fixed form determines $W$
up to isomorphism by its semisimple attachment and hence determines
$1\oplus W|_{G_E}$. Distinct elements of this branch therefore
require distinct normalized primitive eigenforms. These belong
to the standard simultaneous eigenbasis of the full cusp space,
so their eigenvalues occur in $Q_{e,7}$ with sufficient
multiplicity.

After applying $\iota^{-1}$, the two lists of trace differences
are respectively
\[
 \frac{\alpha-(7^{2e}+7^e+1)}{7^e},\qquad
 \frac{a_7-(7^{2e}+1)}{7^e}.
\]
Both characteristic polynomials are monic integral, so all their
roots are integral at $\ell$. Since $\ell\ne7$, the displayed
denominators have valuation zero. Each selected factor has
nonnegative valuation; adding all remaining characteristic
factors only increases the sum. Their product is $N_{e,7}$,
which proves (KT-M.1).

For nonvanishing, equip $V_{e-1}$ with a unitary norm.
The compact Hecke operator is a sum of $57$ unitary operators,
so every scaled eigenvalue has complex absolute value at most
$57\cdot7^{e-1}<7^{2e}+7^e+1$ for $e\geq2$.
The classical bound $|a_7|\leq2\cdot7^e<7^{2e}+1$ gives the
other nonvanishing. Alternatively its weaker elementary Hecke
operator norm also suffices.
The source comparison trace is therefore different from every
admissible trace, and \texttt{OneCharU\_s} applies.
Its finite set is pairwise nonisomorphic by
\texttt{NumberU3US.lean}, lines 57--59, so (KT-M.1) bounds its
valuation sum as asserted.
\end{proof}

\paragraph{A bound for the comparison integer.}
If $D_e=\dim M_{e-1}$ and
$r_e=\dim S_{2e+1}(\Gamma_0(3),\chi_{-3})$, the same norms give
\[
 \log|N_{e,7}|\leq(D_e+r_e)(2e\log7+1)\qquad(e\geq2).
\]
Indeed
$7^{2e}+7^e+1+57\cdot7^{e-1}
 =7^{2e}(1+(64/7)7^{-e}+7^{-2e})$
and the parenthesis is less than $\exp(1)$ for $e\geq2$;
the classical factor satisfies the same estimate.
This is a bound on size, not on smoothness or on the prime
divisors individually. Turning it into an all-prime bound on
$K_{2e}(\mathbf Z)$ additionally requires the exact source
$\mathrm{GD}$ conditions and the inherited Selmer--$K$ comparison
for each prime; those arithmetic conditions must be checked
separately.

\begin{theorem}[KT-QUARTIC-BOUND]
Under the report's carried-source convention, for every even
$e\geq2$,
\[
      |K_{2e}(\mathbf Z)|\mid N_{e,7}^{\,2},\qquad
      \log|K_{2e}(\mathbf Z)|<\frac25 e^4
      \quad(e\geq2\text{ even}).                     \tag{KT-M.2}
\]
The first divisibility is understood with $|N_{e,7}|$ if necessary.
This is an explicit alternative bound; it does not evaluate the
particular existential constant selected inside VK-EXP7.
\end{theorem}
\paragraph{Dependencies.}
The source theorem used is KT-ONECHAR-CLOSED, the closed
proof-term extracted from
\texttt{SmallClassNumbers/CorGKOf.lean}, lines 10--25,
with the \texttt{OneCharU\_s} implication at lines 15--18 and
the published-hypotheses instantiation at lines 20--25.
Its exact implication is in \texttt{DenseThresholdCases/FOneCharU.lean},
lines 75--95. The source chain is read and audited in this report;
it is not a new compilation of the entire Lean project.
The arithmetic inputs are the inherited Selmer--$K$ comparison
VK-PT and exceptional-character vanishing VK-EXCEPTIONS
(\texttt{vandiver\_k.md}, Sections 9.2--9.3), the stipulated
Vandiver verification for odd primes below $2^{31}$ with
Kurihara's zero-equivalence, and the carried Rognes--Weibel
vanishing of the $2$-primary part of $K_{4k}(\mathbf Z)$.
The needed \texttt{GD} assertion is a proved source statement,
\texttt{SmallClassNumbers/ClassNumberThree.lean}, lines 64--86;
its full displayed proof was also read here.
\begin{proof}
For every prime $\ell\geq2^{31}$ with $\ell-1\nmid2e$,
the quoted \texttt{gd} theorem supplies $\mathrm{GD}$ for
$E=\mathbf Q(\sqrt{-3})$, modulus two and trivial character.
Its hypotheses are precisely $e\geq1$, $e$ even,
$\ell\geq5$, and $\ell-1\nmid2e$; no upper bound on $e$
occurs. Apply
KT-TRACE-MULTIPLICITY to the finite set supplied by
$\mathrm{OneCharU\_s}$ and use VK-PT to identify its normalized
length with $v_\ell(|K_{2e}(\mathbf Z)|)$.
If $\ell-1\mid2e$, VK-EXCEPTIONS makes the primary group zero.
All odd primes below $2^{31}$, including seven, contribute zero
by the stated Vandiver theorem and Kurihara's equivalence.
The $2$-primary part is zero by the separate Rognes--Weibel input.
Thus every prime exponent in the group order is bounded by twice
its exponent in $N_{e,7}$, proving the divisibility and
$\log|K_{2e}(\mathbf Z)|\leq2\log|N_{e,7}|$.

The exact dimension formula specializes for positive even $e$ to
\[
D_e=\begin{cases}
e^3/216,&e\equiv0\pmod6,\\
e^3/216+e/9-7/27,&e\equiv2\pmod6,\\
e^3/216+e/9+7/27,&e\equiv4\pmod6.
\end{cases}
\]
The full classical dimension is
$r_e=\lfloor(2e-2)/3\rfloor$.
Hence for $e\geq4$ even,
\[
 S_e:=D_e+r_e\leq e^3/216+7e/9-11/27.
\]
Dividing this right side by $e^3$ gives
$1/216+7/(9e^2)-11/(27e^3)$.
Its derivative is $(33-42e)/(27e^4)<0$ for $e\geq4$,
and its value at four is $3/64$. Therefore $S_e\leq3e^3/64$.
Using the preceding bound for $N_{e,7}$ and $\log7<2$,
\begin{align*}
 \log|K_{2e}(\mathbf Z)|
 &\leq 2S_e(2e\log7+1)\\
 &<\frac3{32}e^3(4e+1)
 \leq\frac{51}{128}e^4<\frac25e^4.
\end{align*}
For $e=2$ both spaces have dimension zero, so the divisibility
already says $|K_4(\mathbf Z)|=1$ and the strict bound follows.
All arithmetic interface conditions have now been supplied by
the listed inputs. Pointwise trace coverage alone would not
give this conclusion: the whole-system injection in
KT-TRACE-MULTIPLICITY is what preserves the necessary
valuation multiplicities.
\end{proof}
